% GENERATED FILE. Edit canonical lessons, then run npm run generate:books.
% canonical-source-hash: fnv1a64:21d58b3fc7695ead
% canonical-lessons: IT-C530-abbonarsi, IT-C530-streaming, IT-C530-videogioco, IT-C530-postare, IT-C530-follower

\chapter{Sullo schermo}
\label{ch:it-sullo-schermo}
\hlchaptermodality{530}

\begin{chapteropening}
\textbf{By the end of this chapter:} \emph{I can talk about streaming, games and social media.}

Complete the last lesson of chapter 530: I can talk about streaming, games and social media.
\end{chapteropening}

\section[abbonarsi]{abbonarsi — to subscribe}

\label{lesson:IT-C530-abbonarsi}



\begin{quote}
Before the new one: say the Italian for soundtrack, then the Italian for trailer.
\end{quote}

\subsection*{You'll want to know: abbonarsi}
\textbf{abbonarsi} — \textquotedblleft{}to subscribe\textquotedblright{}.

\begin{quote}
Mi sono abbonato a una rivista di cinema.
\end{quote}

\begin{quote}
\emph{I subscribed to a film magazine.}
\end{quote}

\emph{Abbonarsi a qualcosa}; it takes \emph{essere}.

\begin{cousinweb}
\emph{abbonarsi} comes from French \emph{abonner} \textquotedblleft{}to subscribe\textquotedblright{}, usually linked to \emph{borne} \textquotedblleft{}a limit\textquotedblright{}; the origin is not fully settled.
\end{cousinweb}

\begin{grammarlens}[title={What you've built}]
One more word for telling, explaining and describing at B1.
\end{grammarlens}

\subsection*{Your turn}
\begin{itemize}
\raggedright
  \item \emph{Say it:} \emph{abbonarsi}
  \item \emph{Say it:} \emph{abbonarsi}, once more
  \item \emph{Say it:} a sentence of your own with \emph{abbonarsi}
  \item \emph{Recall:} say the Italian for soundtrack, then the Italian for trailer, then say \emph{abbonarsi} again
\end{itemize}

\begin{tcolorbox}[breakable,colback=teal!4,colframe=teal!35!black,title={Before you move on}]
What does abbonarsi mean? (\textquotedblleft{}To subscribe\textquotedblright{}.) Say it once more.
\end{tcolorbox}
\section[lo streaming]{lo streaming — streaming}

\label{lesson:IT-C530-streaming}



\begin{quote}
Before the new one: say the Italian for trailer, then the Italian for to subscribe.
\end{quote}

\subsection*{You'll want to know: lo streaming}
\textbf{lo streaming} — \textquotedblleft{}streaming\textquotedblright{}.

\begin{quote}
Guardo i film in streaming sul computer.
\end{quote}

\begin{quote}
\emph{I watch films by streaming on the computer.}
\end{quote}

From English; \emph{in streaming} — streamed online.

\begin{cousinweb}
\emph{streaming} was borrowed from English \emph{streaming}, from \emph{stream}, a Germanic word.
\end{cousinweb}

\begin{grammarlens}[title={What you've built}]
One more word for telling, explaining and describing at B1.
\end{grammarlens}

\subsection*{Your turn}
\begin{itemize}
\raggedright
  \item \emph{Say it:} \emph{lo streaming}
  \item \emph{Say it:} \emph{lo streaming}, once more
  \item \emph{Say it:} a sentence of your own with \emph{lo streaming}
  \item \emph{Recall:} say the Italian for trailer, then the Italian for to subscribe, then say \emph{lo streaming} again
\end{itemize}

\begin{tcolorbox}[breakable,colback=teal!4,colframe=teal!35!black,title={Before you move on}]
What does lo streaming mean? (\textquotedblleft{}Streaming\textquotedblright{}.) Say it once more.
\end{tcolorbox}
\section[il videogioco]{il videogioco — video game}

\label{lesson:IT-C530-videogioco}



\begin{quote}
Before the new one: say the Italian for to subscribe, then the Italian for streaming.
\end{quote}

\subsection*{You'll want to know: il videogioco}
\textbf{il videogioco} — \textquotedblleft{}video game\textquotedblright{}.

\begin{quote}
Mio figlio passa ore con i videogiochi.
\end{quote}

\begin{quote}
\emph{My son spends hours on video games.}
\end{quote}

The plural is \emph{i videogiochi}.

\begin{cousinweb}
\emph{videogioco} is built from Latin \emph{videō} \textquotedblleft{}I see\textquotedblright{} (Lewis \& Short) and \emph{gioco} \textquotedblleft{}game\textquotedblright{}, from Latin \emph{iocus} (EtymDB).
\end{cousinweb}

\begin{grammarlens}[title={What you've built}]
One more word for telling, explaining and describing at B1.
\end{grammarlens}

\subsection*{Your turn}
\begin{itemize}
\raggedright
  \item \emph{Say it:} \emph{il videogioco}
  \item \emph{Say it:} \emph{il videogioco}, once more
  \item \emph{Say it:} a sentence of your own with \emph{il videogioco}
  \item \emph{Recall:} say the Italian for to subscribe, then the Italian for streaming, then say \emph{il videogioco} again
\end{itemize}

\begin{tcolorbox}[breakable,colback=teal!4,colframe=teal!35!black,title={Before you move on}]
What does il videogioco mean? (\textquotedblleft{}Video game\textquotedblright{}.) Say it once more.
\end{tcolorbox}
\section[postare]{postare — to post (online)}

\label{lesson:IT-C530-postare}



\begin{quote}
Before the new one: say the Italian for streaming, then the Italian for video game.
\end{quote}

\subsection*{You'll want to know: postare}
\textbf{postare} — \textquotedblleft{}to post (online)\textquotedblright{}.

\begin{quote}
Ha postato una foto delle vacanze.
\end{quote}

\begin{quote}
\emph{She posted a holiday photo.}
\end{quote}

Informal; more formal \emph{pubblicare}.

\begin{cousinweb}
\emph{postare} is built on English \emph{post} (EtymDB), with the Italian verb ending \emph{-are}.
\end{cousinweb}

\begin{grammarlens}[title={What you've built}]
One more word for telling, explaining and describing at B1.
\end{grammarlens}

\subsection*{Your turn}
\begin{itemize}
\raggedright
  \item \emph{Say it:} \emph{postare}
  \item \emph{Say it:} \emph{postare}, once more
  \item \emph{Say it:} a sentence of your own with \emph{postare}
  \item \emph{Recall:} say the Italian for streaming, then the Italian for video game, then say \emph{postare} again
\end{itemize}

\begin{tcolorbox}[breakable,colback=teal!4,colframe=teal!35!black,title={Before you move on}]
What does postare mean? (\textquotedblleft{}To post (online)\textquotedblright{}.) Say it once more.
\end{tcolorbox}
\section[il follower]{il follower — follower (online)}

\label{lesson:IT-C530-follower}



\begin{quote}
Before the new one: say the Italian for video game, then the Italian for to post.
\end{quote}

\subsection*{You'll want to know: il follower}
\textbf{il follower} — \textquotedblleft{}follower (online)\textquotedblright{}.

\begin{quote}
Ha più di mille follower.
\end{quote}

\begin{quote}
\emph{She has more than a thousand followers.}
\end{quote}

No plural \emph{-s}: \emph{i follower}.

\begin{cousinweb}
\emph{follower} was borrowed from English \emph{follower}, built on \emph{follow}, from Proto-Germanic \emph{fulgijaną} \textquotedblleft{}to follow\textquotedblright{} (EtymDB).
\end{cousinweb}

\begin{grammarlens}[title={What you've built}]
One more word for telling, explaining and describing at B1.
\end{grammarlens}

\subsection*{Your turn}
\begin{itemize}
\raggedright
  \item \emph{Say it:} \emph{il follower}
  \item \emph{Say it:} \emph{il follower}, once more
  \item \emph{Say it:} a sentence of your own with \emph{il follower}
  \item \emph{Recall:} say the Italian for video game, then the Italian for to post, then say \emph{il follower} again
\end{itemize}

\begin{tcolorbox}[breakable,colback=teal!4,colframe=teal!35!black,title={Before you move on}]
What does il follower mean? (\textquotedblleft{}Follower (online)\textquotedblright{}.) Say it once more.
\end{tcolorbox}
